% *************** Zakończenie ***************

%***************************************************************************************
% W tym miejscu znajdują się polecenia odpowiedzialne za tworzenie
% spisu ilustracji, spisu treści oraz streszczenia pracy
%***************************************************************************************

%spis rysunków
\addcontentsline{toc}{chapter}{List of Figures}
\listoffigures
\newpage

%spis tablic
\addcontentsline{toc}{chapter}{List of Tables}
\listoftables
\newpage

%streszczenie
\addcontentsline{toc}{chapter}{Summary}
%\noindent
%{\footnotesize{}\textbf{Wyższa Szkoła Informatyki i Zarządzania z siedzibą w Rzeszowie\\
%\textbf{Kolegium Informatyki Stosowanej}}
%\vspace{30pt}

%\begin{center}
%\textbf{Streszczenie pracy dyplomowej}\\
%\temat
%\end{center}

%\vspace{30pt}
%\noindent
%\textbf{Autor: \autor
%\\Promotor: \promotor
%\\Słowa kluczowe: tutaj umieść słowa kluczowe}
%\vspace{40pt}
%\\Treść streszczenia, czyli kilka zdań dotyczących treści pracy dyplomowej w języku polskim.
%}
%\vspace{80pt}

\noindent
{\footnotesize{}\textbf{University of Information Technology and Management in Rzeszów\\
\textbf{Faculty of Applied Information Technology}}
\vspace{30pt}

\begin{center}
\textbf{Thesis Summary\\}
\temat
\end{center}

\vspace{30pt}
\noindent
\textbf{Author: \autor
\\Supervisor: \promotor
\\Key words: Password Strength, Hybrid Evaluation, Entropy, Pattern Detection, Dictionary Check, CustomTkinter, Python, Usability, Desktop Application, Privacy}
\vspace{40pt}

This thesis presents the design, implementation, and evaluation of a hybrid password strength evaluation tool developed as a Python-based desktop application with a graphical user interface. The tool combines five complementary analysis dimensions---rule-based length scoring, character diversity assessment, Shannon entropy estimation, pattern detection (including reversed keyboard patterns and Levenshtein distance similarity), and dictionary comparison---into a single transparent 0--100 scoring model with qualitative labels and actionable, educational feedback. The application features a dark-themed graphical interface built with CustomTkinter that provides real-time password analysis through a 300\,ms debounced input mechanism, a two-column analysis layout with an animated semicircular gauge, and a toggleable per-dimension score breakdown panel that makes the evaluation process fully transparent to the user. A cryptographically secure password generator using Python's \texttt{secrets} module is also integrated. The tool was evaluated through functional testing with over 20 representative passwords, a comparative analysis against the Norton Password Checker and Kaspersky Password Checker, and a usability self-evaluation based on Nielsen's heuristics. The results demonstrate that the hybrid approach effectively identifies weak passwords that some online checkers miss---particularly reversed keyboard patterns---while maintaining complete user privacy through local-only processing with no data storage or transmission. The thesis documents an iterative development process of 18 systematic enhancements and a user-experience redesign, illustrating how even small security tools benefit from structured quality refinement and transparent, explainable scoring.


% *************** Koniec pliku back.tex ***************
